\chapter{Lista de acrónimos y glosario}

Los términos ingleses sin traducción asentada en español se escriben en cursiva; los ya
incorporados al uso técnico habitual (hardware, software, driver, firmware) se dejan en redonda.
Algunas entradas de esta lista no son siglas, sino términos que conviene precisar; se recogen
aquí por comodidad de consulta.

\paragraph{AOCS} \textit{Attitude and Orbit Control Subsystem} --- Subsistema de control de actitud y órbita.

\paragraph{ADC} \textit{Analog-to-Digital Converter} --- Convertidor analógico-digital.

\paragraph{AMBA} \textit{Advanced Microcontroller Bus Architecture} --- Familia de estándares de interconexión de ARM.

\paragraph{API} \textit{Application Programming Interface} --- Interfaz de programación que un módulo expone a quien lo usa.

\paragraph{ARM} \textit{Advanced RISC Machines} --- Arquitectura de procesador de los núcleos de aplicación del MPSoC.

\paragraph{AXI} \textit{Advanced eXtensible Interface} --- Bus de interconexión de alto rendimiento del estándar ARM AMBA.

\paragraph{BD} \textit{Buffer Descriptor} --- Descriptor de \textit{buffer}: estructura en memoria que define una transferencia de DMA en modo \textit{scatter-gather}.

\paragraph{BER} \textit{Bit Error Rate} --- Tasa de error de bit de un enlace.

\paragraph{BIF} \textit{Boot Image Format} --- Fichero de descripción con el que Vitis empaqueta los componentes de la imagen de arranque.

\paragraph{BRAM} \textit{Block RAM} --- Bloque de memoria RAM integrado en la FPGA.

\paragraph{BSP} \textit{Board Support Package} --- Paquete de soporte de placa.

\paragraph{BOM} \textit{Bill of Materials} --- Lista de materiales de una PCB.

\paragraph{CAN} \textit{Controller Area Network} --- Bus de comunicaciones serie multipunto.

\paragraph{CDHS} \textit{Command and Data Handling System} --- Subsistema de mando y gestión de datos.

\paragraph{CMRR} \textit{Common Mode Rejection Ratio} --- Razón de rechazo al modo común.

\paragraph{COTS} \textit{Commercial Off-The-Shelf} --- Componente comercial de catálogo, no cualificado específicamente para espacio.

\paragraph{CRC} \textit{Cyclic Redundancy Check} --- Código de redundancia cíclica para la detección de errores de transmisión.

\paragraph{CPU} \textit{Central Processing Unit} --- Unidad central de proceso.

\paragraph{DDR} \textit{Double Data Rate} --- Memoria dinámica principal del sistema de procesamiento.

\paragraph{DMA} \textit{Direct Memory Access} --- Acceso directo a memoria sin intervención de la CPU.

\paragraph{DSP} \textit{Digital Signal Processor} --- Bloque de procesamiento digital de señal.

\paragraph{DTB} \textit{Device Tree Binary} --- Árbol de dispositivos en formato binario para sistemas Linux/U-Boot.

\paragraph{EMC} \textit{Electromagnetic Compatibility} --- Compatibilidad electromagnética.

\paragraph{EMI} \textit{Electromagnetic Interference} --- Interferencia electromagnética.

\paragraph{ESD} \textit{Electrostatic Discharge} --- Descarga electrostática.

\paragraph{ESA} \textit{European Space Agency} --- Agencia Espacial Europea.

\paragraph{ETSIT} Escuela Técnica Superior de Ingenieros de Telecomunicación.

\paragraph{FIFO} \textit{First In, First Out} --- Estructura de datos de cola.

\paragraph{FMC} \textit{FPGA Mezzanine Card} --- Tarjeta de expansión para FPGA según estándar VITA 57.1.

\paragraph{FPGA} \textit{Field Programmable Gate Array} --- Matriz lógica programable en campo.

\paragraph{FSM} \textit{Finite State Machine} --- Máquina de estados finitos.

\paragraph{FSBL} \textit{First Stage Boot Loader} --- Cargador de arranque de primera etapa.

\paragraph{GIC} \textit{Generic Interrupt Controller} --- Controlador de interrupciones del procesador ARM.

\paragraph{GPIO} \textit{General Purpose Input/Output} --- Puerto de entrada/salida de propósito general.

\paragraph{GUI} \textit{Graphical User Interface} --- Interfaz gráfica de usuario.

\paragraph{HDL} \textit{Hardware Description Language} --- Lenguaje de descripción de hardware.

\paragraph{HP} \textit{High Performance} --- Puertos AXI de alto rendimiento entre el sistema de procesamiento y la lógica programable del MPSoC.

\paragraph{HPC} \textit{High Pin Count} --- Variante del conector FMC con el mayor número de pines.

\paragraph{IDE} \textit{Integrated Development Environment} --- Entorno integrado de desarrollo.

\paragraph{INTC} \textit{Interrupt Controller} --- Controlador de interrupciones (IP de la lógica programable).

\paragraph{IRQ} \textit{Interrupt Request} --- Petición de interrupción.

\paragraph{ISR} \textit{Interrupt Service Routine} --- Rutina de servicio de interrupción.

\paragraph{LDO} \textit{Low-Dropout Regulator} --- Regulador lineal de baja caída de tensión.

\paragraph{LEO} \textit{Low Earth Orbit} --- Órbita terrestre baja.

\paragraph{LINCE} \textit{Línea de INdustrialización de Cargas de pago y plataformas Espaciales}.

\paragraph{LVDS} \textit{Low-Voltage Differential Signaling} --- Señalización diferencial de bajo voltaje.

\paragraph{LUT} \textit{Look-Up Table} --- Tabla de búsqueda, bloque básico de una FPGA.

\paragraph{LSB} \textit{Least Significant Bit} --- Bit menos significativo; en el transceptor, uno de los dos órdenes de transmisión.

\paragraph{MCDMA} \textit{Multichannel Direct Memory Access} --- Controlador de DMA multicanal de AMD (IP \texttt{axi\_mcdma}, PG288).

\paragraph{MMIO} \textit{Memory-Mapped I/O} --- Entrada/salida con mapeado en memoria.

\paragraph{MM2S} \textit{Memory-Mapped to Stream} --- Canal de un DMA que lee de memoria y produce un flujo AXI-Stream (camino de transmisión).

\paragraph{MMU} \textit{Memory Management Unit} --- Unidad de gestión de memoria; define, entre otras cosas, qué regiones son cacheables.

\paragraph{MPSoC} \textit{Multiprocessor System-on-Chip} --- Sistema multiprocesador en chip.

\paragraph{MSB} \textit{Most Significant Bit} --- Bit más significativo.

\paragraph{NCO} \textit{Numerically Controlled Oscillator} --- Oscilador controlado numéricamente.

\paragraph{OBC} \textit{On-Board Computer} --- Ordenador de a bordo.

\paragraph{PCB} \textit{Printed Circuit Board} --- Placa de circuito impreso.

\paragraph{PERTE} Proyecto Estratégico para la Recuperación y Transformación Económica --- Instrumento de financiación pública del Gobierno de España.

\paragraph{PL} \textit{Programmable Logic} --- Lógica programable (FPGA del MPSoC).

\paragraph{PMUFW} \textit{Power Management Unit Firmware} --- Firmware de la unidad de gestión de potencia.

\paragraph{PS} \textit{Processing System} --- Sistema de procesamiento (ARM del MPSoC).

\paragraph{PWM} \textit{Pulse Width Modulation} --- Modulación por ancho de pulso.

\paragraph{PYMES} Pequeñas y medianas empresas.

\paragraph{QEMU} \textit{Quick Emulator} --- Emulador de máquina que ejecuta software sin el hardware físico.

\paragraph{RSB} \textit{RTEMS Source Builder} --- Herramienta que compila el \textit{toolchain} y los BSP de RTEMS desde código fuente.

\paragraph{RTEMS} \textit{Real-Time Executive for Multiprocessor Systems} --- RTOS de código abierto para sistemas empotrados.

\paragraph{RTL} \textit{Register Transfer Level} --- Nivel de descripción del hardware mediante transferencias entre registros.

\paragraph{RTOS} \textit{Real-Time Operating System} --- Sistema operativo de tiempo real.

\paragraph{S2MM} \textit{Stream to Memory-Mapped} --- Canal de un DMA que consume un flujo AXI-Stream y lo escribe en memoria (camino de recepción).

\paragraph{SAR} \textit{Successive Approximation Register} --- Arquitectura de convertidor analógico-digital por aproximaciones sucesivas.

\paragraph{SEU} \textit{Single Event Upset} --- Cambio de estado de un biestable o de una celda de memoria provocado por el impacto de una partícula.

\paragraph{SG} \textit{Scatter-Gather} --- Modo de operación de un DMA en el que las transferencias se describen mediante una lista enlazada de descriptores en memoria.

\paragraph{SLO} Pin de los transceptores LTC2865 que habilita el limitador de velocidad de flanco, activo a nivel bajo.

\paragraph{SMD} \textit{Surface Mount Device} --- Componente de montaje superficial.

\paragraph{SMP} \textit{Symmetric Multiprocessing} --- Multiprocesamiento simétrico sobre varios núcleos.

\paragraph{SPI} \textit{Serial Peripheral Interface} --- Interfaz serie para periféricos.

\paragraph{SPW} Abreviatura de SpaceWire empleada en los mapas de señales y los ficheros de restricciones.

\paragraph{TCL} \textit{Tool Command Language} --- Lenguaje de scripting utilizado en Vivado.

\paragraph{TDEST} \textit{Transfer Destination} --- Señal del canal AXI-Stream que identifica el canal de destino de cada paquete.

\paragraph{TFM} Trabajo Fin de Máster.

\paragraph{TLAST} \textit{Transfer Last} --- Señal del canal AXI-Stream que marca el último byte de un paquete.

\paragraph{TMR} \textit{Triple Modular Redundancy} --- Triplicación de la lógica con votadores para tolerar los fallos provocados por la radiación.

\paragraph{TREADY} \textit{Transfer Ready} --- Señal del canal AXI-Stream con la que el receptor indica que puede aceptar un dato.

\paragraph{TVALID} \textit{Transfer Valid} --- Señal del canal AXI-Stream con la que el emisor indica que el dato presentado es válido.

\paragraph{TVS} \textit{Transient Voltage Suppressor} --- Supresor de sobretensiones transitorias.

\paragraph{UART} \textit{Universal Asynchronous Receiver-Transmitter} --- Transceptor serie asíncrono universal.

\paragraph{UPM} Universidad Politécnica de Madrid.

\paragraph{USB} \textit{Universal Serial Bus} --- Bus serie empleado aquí para la consola de la ZCU102.

\paragraph{VADJ} Carril de alimentación ajustable con el que la ZCU102 alimenta los bancos de entrada/salida del conector FMC; en este trabajo, a 1,8~V.

\paragraph{VHDL} \textit{Very High Speed Integrated Circuit Hardware Description Language} --- Lenguaje de descripción de hardware.

\paragraph{WNS} \textit{Worst Negative Slack} --- Peor margen de temporización de un diseño implementado; positivo indica que se cumplen las restricciones.

\paragraph{XSA} \textit{Xilinx Support Archive} --- Archivo de especificación hardware exportado por Vivado para Vitis.
